\documentclass[11pt,a4paper]{article}
\usepackage[T1]{fontenc}
\usepackage[utf8]{inputenc}
\usepackage[margin=2.5cm]{geometry}
\usepackage{enumitem}
\setlist{itemsep=2pt}

\title{Two changes proposed to the PET CAEN DAQ}
\author{J.\,J.~Gómez-Cadenas and Claude (Anthropic)}
\date{\today}

\begin{document}
\maketitle

The CRYSP bench calibrates the Citiroc~1A on the A5202, read through the
DT5215 with pet-caen-daq, before taking data with BGO and CsI(Tl) crystals.
Reading the DAQ next to FERSlib 5.0.1 we found two things that matter for
these measurements, and opened one pull request for each. Both pass the
complete suite that runs without hardware (Go plain, HDF5 and simulator
integration tests; frontend format, lint, type check and unit tests).

\begin{enumerate}
\item {\bf \#73: write the \texttt{CitirocConfig} register on every
  configuration.} With automatic ASIC configuration the FPGA builds the
  slow-control stream of each Citiroc from board registers, and the chip
  modes come from \texttt{CitirocConfig} (\texttt{0x01000108}). FERSlib
  composes this register on every configuration; the DAQ writes it only
  during a threshold staircase, so a data run keeps whatever value the board
  held before. As a consequence \texttt{FastShaperInput} and
  \texttt{HV\_Adjust\_Range} are validated and audited as applied but do not
  reach the chip. The pull request composes the register as FERSlib does
  (\texttt{0x00070f21} with the JANUS defaults, \texttt{0x00070fa1} for the
  production sample) and adds it to the configuration plan.
\item {\bf \#74: add \texttt{PeakDetectorMode}} (PEAK\_STRETCH or
  TRACK\&HOLD), built on \#73. FERSlib reads it, the JANUS GUI does not list
  it, and the DAQ catalogue rejects it, so today every run is taken in peak
  stretch. Omitted, it means PEAK\_STRETCH and gives the register of \#73,
  so existing configurations are unchanged. We will calibrate and take data
  in peak stretch; the setting is for a short check later in track and hold,
  three periodic-trigger runs whose spectra are sensitive to the shape of
  the shaped pulse and are compared with our simulation.
\end{enumerate}

A third pull request (\#75, full-resolution spectra in the hold-delay
scan) is closed: we no longer need it, and its split by the discriminator
flag, which is the latched hit register and carries no timing, would not
select the channel that triggered.

\paragraph{Questions.} The DAQ does not expose \texttt{EnableQdiscrLatch}
and \texttt{EnableChannelTrgout}; \#73 uses the FERSlib defaults (both 1).
Is that the value wanted in production? And was the absence of
\texttt{CitirocConfig} from the configuration plan deliberate, for a
hardware reason we may have missed?

\end{document}
